\uuid{Mfbe}
\exo7id{6876}
\titre{exo7 6876}
\auteur{gammella}
\organisation{exo7}
\datecreate{2012-05-29}
\isIndication{false}
\isCorrection{true}
\chapitre{Analyse vectorielle}
\sousChapitre{Forme différentielle, champ de vecteurs, circulation}
\cpge{mpsi-mp2i.fonctions-deux-variables.b, pcsi.fonctions-deux-variables.b, ptsi.fonctions-deux-variables.b, bcpst1.fonctions-deux-variables.c, mp-mpi.calcul-differentiel-optimisation.b, pc.fonctions-plusieurs-variables.c, psi.fonctions-plusieurs-variables.c}
\module{Géométrie}
\niveau{L2}
\difficulte{}

\contenu{
\question{
On considère le champ vectoriel $\vec{V}(x,y)=(1+2xy,x^3-3)$. 
Ce champ est-il un champ de gradient ?
}
\reponse{
Au champ $\vec{V}(x,y)$ est associée la forme
$$\omega=(1+2xy) dx + (x^3-3) dy.$$
Cette forme n'est pas exacte puisque
$ \frac{\partial (1+2xy)}{\partial y}\not=  \frac{\partial(x^3-3)}{\partial x}$. Il s'ensuit
que $\vec{V(x,y)}$ n'est pas un champ de gradient.
}
}
